\documentclass[a4paper,landscape,8pt]{extarticle}
\usepackage[landscape,margin=0.7cm,top=0.5cm,bottom=0.7cm]{geometry}
\usepackage[utf8]{inputenc}
\usepackage[T1]{fontenc}
\usepackage[ngerman]{babel}
\usepackage{helvet}
\renewcommand{\familydefault}{\sfdefault}
\usepackage{xcolor}
\usepackage{tikz}
\usetikzlibrary{shapes, arrows.meta, positioning}
\usepackage{enumitem}
\usepackage{multicol}
\usepackage{array}
\usepackage{fancyhdr}
\usepackage{amssymb}

% Farben
\definecolor{findmygreen}{HTML}{34C759}
\definecolor{appleblue}{HTML}{007AFF}
\definecolor{lightgray}{HTML}{F5F5F5}
\definecolor{bordergray}{HTML}{BBBBBB}
\definecolor{warnred}{HTML}{C62828}

\pagestyle{fancy}
\fancyhf{}
\renewcommand{\headrulewidth}{0pt}
\fancyfoot[C]{\footnotesize Seite \thepage{} von 3}

\setlength{\parindent}{0pt}
\setlength{\parskip}{2pt}

% Kompakte Listen
\setlist{nosep, leftmargin=1em, itemsep=0pt, topsep=1pt}

\begin{document}

% ============================================================================
% SEITE 1: INFORMATIONSBLATT
% ============================================================================

\noindent{\Large\bfseries ``Wo ist?'' -- Geräte finden \& schützen} \hfill
{\footnotesize\textbf{Lernziele:} ``Wo ist?'' verstehen | Einrichten | Bei Verlust handeln}

\vspace{1mm}
\hrule
\vspace{1mm}

% Drei-Spalten-Layout
\begin{minipage}[t]{0.32\linewidth}
\begin{center}
\colorbox{findmygreen!15}{%
\begin{minipage}{0.95\linewidth}
\vspace{2mm}
\begin{center}
{\Large\bfseries\textcolor{findmygreen}{1. Was ist ``Wo ist?''}}\\[1pt]
{\footnotesize\itshape Alle Geräte auf einer Karte}
\end{center}
\vspace{2mm}
\end{minipage}%
}
\end{center}

\vspace{1mm}

\textbf{Dein digitaler GPS-Tracker}\\
``Wo ist?'' zeigt dir, wo \textbf{alle} deine Apple-Geräte sind -- auf einer Karte.

\textbf{Was kann ``Wo ist?''?}
\begin{itemize}
\item iPhone, iPad, Mac, AirPods orten
\item Ton auf dem Gerät abspielen
\item Gerät als \textbf{verloren markieren}
\item Gerät aus der Ferne \textbf{löschen}
\end{itemize}

\textbf{Analogie: Wie ein Schlüsselfinder}
\begin{itemize}
\item Kennst du ``Tile'' oder ``AirTag''?
\item ``Wo ist?'' ist wie ein eingebauter Finder
\item Funktioniert sogar wenn iPhone \textbf{aus} ist!
\end{itemize}

\vspace{2mm}
\fcolorbox{findmygreen}{findmygreen!5}{%
\begin{minipage}{0.92\linewidth}
\footnotesize
\textbf{Wichtig:} Funktioniert nur, wenn \textbf{vorher aktiviert}! Darum jetzt einrichten.
\end{minipage}%
}

\end{minipage}%
\hfill%
\begin{minipage}[t]{0.32\linewidth}
\begin{center}
\colorbox{appleblue!15}{%
\begin{minipage}{0.95\linewidth}
\vspace{2mm}
\begin{center}
{\Large\bfseries\textcolor{appleblue}{2. ``Wo ist?'' einrichten}}\\[1pt]
{\footnotesize\itshape Drei wichtige Schalter}
\end{center}
\vspace{2mm}
\end{minipage}%
}
\end{center}

\vspace{1mm}

\textbf{So aktivierst du es:}
\begin{enumerate}
\item Öffne \textbf{Einstellungen}
\item Tippe auf \textbf{[Dein Name]} (ganz oben)
\item Tippe auf \textbf{``Wo ist?''}
\item Tippe auf \textbf{``Mein iPhone suchen''}
\end{enumerate}

\textbf{Diese Schalter aktivieren:}

\fcolorbox{bordergray}{lightgray}{%
\begin{minipage}{0.92\linewidth}
\footnotesize
\begin{itemize}[leftmargin=0.8em]
\item[$\square$] \textbf{Mein iPhone suchen} = Grundfunktion
\item[$\square$] \textbf{Netzwerk ``Wo ist?''} = Auch wenn offline
\item[$\square$] \textbf{Letzten Standort senden} = Bei leerem Akku
\end{itemize}
\end{minipage}%
}

\vspace{2mm}

\textbf{Prüfe jetzt:} Geh in die Einstellungen und schau nach, ob alle drei Schalter \textbf{grün} sind!

\end{minipage}%
\hfill%
\begin{minipage}[t]{0.32\linewidth}
\begin{center}
\colorbox{warnred!15}{%
\begin{minipage}{0.95\linewidth}
\vspace{2mm}
\begin{center}
{\Large\bfseries\textcolor{warnred}{3. Gerät verloren?}}\\[1pt]
{\footnotesize\itshape Was du sofort tun kannst}
\end{center}
\vspace{2mm}
\end{minipage}%
}
\end{center}

\vspace{1mm}

\textbf{Schritt-für-Schritt:}
\begin{enumerate}
\item \textbf{Anderes Apple-Gerät?}\\
      $\rightarrow$ Öffne App ``Wo ist?''
\item \textbf{Kein Gerät?}\\
      $\rightarrow$ Browser: \textbf{icloud.com/find}
\item \textbf{Auf Karte suchen}\\
      $\rightarrow$ Gerät antippen
\item \textbf{Ton abspielen}\\
      $\rightarrow$ Wenn in der Nähe
\item \textbf{Als verloren markieren}\\
      $\rightarrow$ Sperrt + zeigt Nachricht
\item \textbf{Bei Diebstahl: Löschen}\\
      $\rightarrow$ Alle Daten weg
\end{enumerate}

\vspace{2mm}
\fcolorbox{warnred}{warnred!5}{%
\begin{minipage}{0.92\linewidth}
\footnotesize
\textbf{Tipp:} Zeige einer \textbf{Vertrauensperson}, wie icloud.com/find funktioniert!
\end{minipage}%
}

\end{minipage}

\vspace{1mm}

% Zusammenfassung
\begin{center}
\fcolorbox{bordergray}{lightgray}{%
\begin{minipage}{0.98\linewidth}
\centering
\vspace{1.5mm}
\textbf{Zusammenfassung:}
\textcolor{findmygreen}{\textbf{``Wo ist?''}} = Geräte auf Karte finden ~$|$~
\textbf{Vorher aktivieren} (3 Schalter) ~$|$~
Bei Verlust: \textbf{Ton abspielen} oder \textbf{als verloren markieren} ~$|$~
Im Browser: \textbf{icloud.com/find}
\vspace{1.5mm}
\end{minipage}%
}
\end{center}

\newpage

% ============================================================================
% SEITE 2: ÜBUNGEN
% ============================================================================

\begin{center}
{\LARGE\bfseries Übungen: ``Wo ist?'' verstehen und anwenden}\\[3pt]
{\normalsize Prüfe dein Wissen und übe am eigenen Gerät}
\end{center}

\vspace{2mm}
\hrule
\vspace{3mm}

\begin{multicols}{2}

\textbf{\large Teil A: Richtig oder Falsch?}

\vspace{2mm}

\renewcommand{\arraystretch}{1.6}
\begin{tabular}{@{}p{8.5cm}cc@{}}
& R & F \\
1. ``Wo ist?'' funktioniert auch ohne vorherige Aktivierung. & $\square$ & $\square$ \\
2. Man kann einen Ton auf dem verlorenen iPhone abspielen. & $\square$ & $\square$ \\
3. ``Wo ist?'' funktioniert nur, wenn das iPhone eingeschaltet ist. & $\square$ & $\square$ \\
4. Ich kann mein iPhone im Browser unter icloud.com finden. & $\square$ & $\square$ \\
5. Wenn ich mein iPhone lösche, sind alle Daten weg. & $\square$ & $\square$ \\
6. Die Aktivierungssperre schützt vor Diebstahl. & $\square$ & $\square$ \\
\end{tabular}
\renewcommand{\arraystretch}{1}

{\footnotesize\textit{Lösungen: F, R, F (mit ``Netzwerk Wo ist?''), R, R, R}}

\vspace{6mm}

\textbf{\large Teil B: Was machst du?}

Welche ``Wo ist?''-Funktion nutzt du?

\vspace{3mm}

\renewcommand{\arraystretch}{1.8}
\begin{tabular}{@{}p{7.5cm}p{4cm}@{}}
1. Du hast dein iPhone irgendwo in der Wohnung verlegt: & \dotfill \\
2. Dein iPhone wurde gestohlen: & \dotfill \\
3. Du hast dein iPhone im Zug liegen lassen: & \dotfill \\
4. Du willst deiner Familie zeigen, wo du bist: & \dotfill \\
\end{tabular}
\renewcommand{\arraystretch}{1}

{\footnotesize\textit{Lösungen: Ton abspielen, Löschen (+ Polizei), Als verloren markieren, Standort teilen}}

\columnbreak

\textbf{\large Teil C: Praktische Übung}

Mache das jetzt an deinem iPhone:

\vspace{3mm}

\fcolorbox{findmygreen}{findmygreen!8}{%
\begin{minipage}{0.95\linewidth}
\vspace{2mm}
\textbf{Übung 1: Einstellungen prüfen}
\begin{enumerate}[leftmargin=1.2em, itemsep=2pt]
\item Gehe zu Einstellungen $\rightarrow$ [Dein Name]
\item Tippe auf ``Wo ist?''
\item Sind alle \textbf{drei} Schalter grün?
\item[$\square$] Ja, alles aktiviert
\item[$\square$] Nein $\rightarrow$ jetzt aktivieren!
\end{enumerate}
\vspace{2mm}
\end{minipage}}

\vspace{4mm}

\fcolorbox{appleblue}{appleblue!8}{%
\begin{minipage}{0.95\linewidth}
\vspace{2mm}
\textbf{Übung 2: App ausprobieren}
\begin{enumerate}[leftmargin=1.2em, itemsep=2pt]
\item Öffne die App ``Wo ist?''
\item Siehst du dein iPhone auf der Karte?
\item Tippe auf ``Ton abspielen''
\item[$\square$] Hat funktioniert!
\item[$\square$] Brauche Hilfe
\end{enumerate}
\vspace{2mm}
\end{minipage}}

\vspace{4mm}

\fcolorbox{bordergray}{lightgray}{%
\begin{minipage}{0.95\linewidth}
\vspace{2mm}
\textbf{Übung 3: Browser testen}

Zu Hause am Computer ausprobieren:
\begin{enumerate}[leftmargin=1.2em, itemsep=2pt]
\item Gehe zu \textbf{icloud.com/find}
\item Melde dich mit deiner Apple-ID an
\item Findest du dein iPhone?
\end{enumerate}
\vspace{2mm}
\end{minipage}}

\end{multicols}

\newpage

% ============================================================================
% SEITE 3: CHECKLISTE + NOTFALL-KARTE
% ============================================================================

\begin{center}
{\LARGE\bfseries Checkliste \& Notfall-Karte}\\[3pt]
{\normalsize Was du verstanden haben solltest + Notfall-Anleitung zum Aufbewahren}
\end{center}

\vspace{3mm}
\hrule
\vspace{4mm}

% Zwei Spalten: Checkliste links, Notfall-Karte rechts
\begin{minipage}[t]{0.48\linewidth}
\begin{center}
\fcolorbox{findmygreen}{findmygreen!10}{%
\begin{minipage}{0.95\linewidth}
\vspace{3mm}
\begin{center}
{\large\bfseries\textcolor{findmygreen}{Checkliste: Das habe ich verstanden}}
\end{center}
\vspace{3mm}

\textbf{Grundlagen:}
\begin{itemize}[label=$\square$, itemsep=4pt]
\item Ich weiß, was ``Wo ist?'' macht
\item Ich verstehe, warum ich es \textbf{vorher} aktivieren muss
\item Ich kenne die drei wichtigen Schalter
\end{itemize}

\vspace{2mm}

\textbf{Bedienung:}
\begin{itemize}[label=$\square$, itemsep=4pt]
\item Ich kann die ``Wo ist?''-App öffnen
\item Ich kann mein Gerät auf der Karte finden
\item Ich kann einen Ton abspielen lassen
\end{itemize}

\vspace{2mm}

\textbf{Notfall:}
\begin{itemize}[label=$\square$, itemsep=4pt]
\item Ich weiß, was \textbf{icloud.com/find} ist
\item Ich kenne den Unterschied: ``verloren'' vs. ``löschen''
\item Ich habe einer Vertrauensperson gezeigt, wie es geht
\end{itemize}

\vspace{3mm}
\end{minipage}%
}
\end{center}

\vspace{4mm}

\begin{center}
\fcolorbox{bordergray}{lightgray}{%
\begin{minipage}{0.95\linewidth}
\vspace{3mm}
\begin{center}
{\bfseries Meine ``Wo ist?''-Einstellungen}
\end{center}
\vspace{2mm}

\renewcommand{\arraystretch}{1.5}
\begin{tabular}{@{}p{6cm}l@{}}
\textbf{``Mein iPhone suchen'' aktiviert?} & $\square$ Ja ~~ $\square$ Nein \\
\textbf{``Netzwerk Wo ist?'' aktiviert?} & $\square$ Ja ~~ $\square$ Nein \\
\textbf{``Letzten Standort senden'' aktiviert?} & $\square$ Ja ~~ $\square$ Nein \\
\end{tabular}
\renewcommand{\arraystretch}{1}
\vspace{3mm}
\end{minipage}%
}
\end{center}

\end{minipage}%
\hfill%
\begin{minipage}[t]{0.48\linewidth}

% Notfall-Karte zum Ausschneiden
\fcolorbox{warnred}{white}{%
\begin{minipage}{0.97\linewidth}
\vspace{3mm}
\begin{center}
{\large\bfseries\textcolor{warnred}{NOTFALL-KARTE}}\\[2pt]
{\footnotesize Ausschneiden und aufbewahren!}
\end{center}

\vspace{2mm}
\hrule
\vspace{3mm}

{\bfseries\textcolor{warnred}{iPhone verloren oder gestohlen?}}

\vspace{2mm}

\textbf{Schritt 1:} Anderes Apple-Gerät nutzen
\begin{itemize}[leftmargin=1.2em, itemsep=1pt]
\item App ``Wo ist?'' öffnen
\item Auf Karte schauen
\end{itemize}

\textbf{Schritt 2:} Kein Apple-Gerät? Computer nutzen:
\begin{itemize}[leftmargin=1.2em, itemsep=1pt]
\item Browser öffnen (Safari, Chrome...)
\item Gehe zu: \textbf{icloud.com/find}
\item Mit Apple-ID anmelden
\end{itemize}

\textbf{Schritt 3:} Gerät gefunden?
\begin{itemize}[leftmargin=1.2em, itemsep=1pt]
\item \textbf{In der Nähe:} ``Ton abspielen''
\item \textbf{Woanders:} ``Als verloren markieren''
\item \textbf{Gestohlen:} ``iPhone löschen'' + Polizei
\end{itemize}

\vspace{2mm}
\hrule
\vspace{2mm}

\textbf{Meine Daten:}\\[2mm]
\renewcommand{\arraystretch}{1.4}
\begin{tabular}{@{}p{3cm}p{5cm}@{}}
Apple-ID: & \dotfill \\
Vertrauensperson: & \dotfill \\
Telefon: & \dotfill \\
\end{tabular}
\renewcommand{\arraystretch}{1}

\vspace{3mm}
\end{minipage}%
}

\vspace{4mm}

\begin{center}
\fcolorbox{appleblue}{appleblue!5}{%
\begin{minipage}{0.95\linewidth}
\vspace{2mm}
\centering
\footnotesize
\textbf{Tipp:} Zeige diese Karte deiner Familie oder einer Vertrauensperson.\\
So kann dir im Notfall jemand helfen!
\vspace{2mm}
\end{minipage}%
}
\end{center}

\end{minipage}

\vfill

\begin{center}
\footnotesize\textit{Stand: \today{} -- ``Wo ist?'' nur nützlich wenn VORHER aktiviert. Heute einrichten!}
\end{center}

\end{document}
